% =====================================================================
%  ARRIVAL Protocol — Full Research Report for Zenodo
%  DOI: 10.5281/zenodo.18893515
%  License: CC BY-NC 4.0 (text), AGPL-3.0-or-later (code)
%  Author: Mefodiy Kelevra — ORCID 0009-0003-4153-392X
% =====================================================================
\documentclass[11pt,a4paper]{article}

% ── Geometry ──────────────────────────────────────────────────────────
\usepackage[margin=2.4cm,top=2.8cm,bottom=2.8cm]{geometry}

% ── Fonts & Encoding ─────────────────────────────────────────────────
\usepackage[T1]{fontenc}
\usepackage[utf8]{inputenc}
\usepackage{lmodern}
\usepackage{microtype}

% ── Math ─────────────────────────────────────────────────────────────
\usepackage{amsmath,amssymb,amsthm}

% ── Tables & Figures ─────────────────────────────────────────────────
\usepackage{booktabs}
\usepackage{tabularx}
\usepackage{multirow}
\usepackage{graphicx}
\usepackage{float}

% ── Colours & Boxes ──────────────────────────────────────────────────
\usepackage[dvipsnames]{xcolor}
\usepackage[most]{tcolorbox}

% ── Hyperlinks ───────────────────────────────────────────────────────
\usepackage[colorlinks=true,linkcolor=NavyBlue,citecolor=OliveGreen,
            urlcolor=BrickRed]{hyperref}
\usepackage{url}

% ── Bibliography ─────────────────────────────────────────────────────
\usepackage[numbers,sort&compress]{natbib}

% ── Section formatting ───────────────────────────────────────────────
\usepackage{titlesec}
\titleformat{\section}{\Large\bfseries}{\thesection.}{0.5em}{}
\titleformat{\subsection}{\large\bfseries}{\thesubsection}{0.5em}{}
\titleformat{\subsubsection}{\normalsize\bfseries}{\thesubsubsection}{0.5em}{}

% ── Theorem-like environments (tcolorbox style) ─────────────────────
\newtcolorbox{theorembox}[1][]{
  colback=blue!3, colframe=NavyBlue!70!black,
  fonttitle=\bfseries, title={#1},
  boxrule=0.6pt, arc=2pt, left=6pt, right=6pt, top=4pt, bottom=4pt
}
\newtcolorbox{definitionbox}[1][]{
  colback=green!3, colframe=OliveGreen!70!black,
  fonttitle=\bfseries, title={#1},
  boxrule=0.6pt, arc=2pt, left=6pt, right=6pt, top=4pt, bottom=4pt
}
\newtcolorbox{resultbox}[1][]{
  colback=red!3, colframe=BrickRed!60!black,
  fonttitle=\bfseries, title={#1},
  boxrule=0.6pt, arc=2pt, left=6pt, right=6pt, top=4pt, bottom=4pt
}

% ── Misc ─────────────────────────────────────────────────────────────
\usepackage{enumitem}
\setlist{nosep,leftmargin=*}
\usepackage{xspace}
\newcommand{\arrival}{\textsc{Arrival}\xspace}
\newcommand{\cgd}{\textsc{CGD}\xspace}
\newcommand{\care}{\textsc{Care}\xspace}
\newcommand{\pp}{\,pp\xspace}

% ── Header / Footer ──────────────────────────────────────────────────
\usepackage{fancyhdr}
\pagestyle{fancy}
\fancyhf{}
\fancyhead[L]{\small\textit{ARRIVAL Protocol — Research Report}}
\fancyhead[R]{\small\textit{DOI: 10.5281/zenodo.18893515}}
\fancyfoot[C]{\thepage}
\renewcommand{\headrulewidth}{0.4pt}

% =====================================================================
\begin{document}

% ── Title Page ───────────────────────────────────────────────────────
\begin{titlepage}
\centering
\vspace*{2cm}

{\Huge\bfseries ARRIVAL Protocol}\\[0.6cm]
{\Large Structured Semantic Coordination Enables\\
Sustainable Cross-Architecture AI Cooperation}\\[2.5cm]

{\large Mefodiy Kelevra}\\[0.3cm]
{\small ORCID:
  \href{https://orcid.org/0009-0003-4153-392X}{0009-0003-4153-392X}}\\[2cm]

{\normalsize March 2026}\\[0.5cm]
{\small DOI: \href{https://doi.org/10.5281/zenodo.18893515}%
  {10.5281/zenodo.18893515}}\\[2cm]

\rule{0.6\textwidth}{0.4pt}\\[0.8cm]

{\footnotesize
  \textbf{License} — Text: CC BY-NC 4.0 \enspace|\enspace
  Code: AGPL-3.0-or-later\\[0.4cm]
  Repository:
  \url{https://github.com/DreamOS-Network/Arrival-Protocol}\\[0.2cm]
  Mathematical foundation: MEANING-CRDT v1.1
  (DOI: \href{https://doi.org/10.5281/zenodo.18702383}%
  {10.5281/zenodo.18702383})
}

\vfill
\end{titlepage}

% ── Table of Contents ────────────────────────────────────────────────
\tableofcontents
\newpage

% =====================================================================
\section{Abstract}
% =====================================================================

We present the \arrival Protocol (Atomic Reasoning via Rival Validation
and Adversarial Logic), a structured communication framework enabling
AI-to-AI coordination across heterogeneous large language model (LLM)
architectures without fine-tuning, shared weights, or prior joint
training.  The protocol employs 66 semantic @-atoms from
DEUS.PROTOCOL~v0.5, injected via system prompts, to establish a shared
coordination vocabulary.

Across 2,200+ experiments organized into 23~phases involving
17~distinct LLM architectures, the principal findings are:

\begin{itemize}
  \item \textbf{GPQA Diamond (N=40, Phase~16):}
        65.0\% accuracy with a homogeneous 5-agent Qwen3-235B ensemble
        vs.\ 52.5\% majority voting (+12.5\pp, McNemar $p=0.006$).
  \item \textbf{Full GPQA Diamond (N=198, Phase~20):}
        \arrival MV 67.7\% vs.\ Solo CoT 62.1\% (+5.6\pp,
        $p=0.233$, ns).
  \item \textbf{Frontier models (Phase~21):}
        Non-Debate MV 82.3\% vs.\ \arrival MV 78.8\% ($-3.5$\pp).
        Anchoring identified as primary degradation mechanism.
  \item \textbf{CGD (Phase~22, N=198):}
        \textbf{86.4\%} accuracy (171/198, 95\% CI [80.9\%, 90.5\%]);
        +7.6\pp over \arrival MV ($p=0.009$);
        61.6\% questions unanimous (95.9\% acc.).
  \item \textbf{CGD-7 (Phase~23, N=198):}
        86.9\% with 7 cross-vendor models; best solo Grok~4.1 87.4\%
        outperforms ensemble — quality $>$ quantity.
  \item \textbf{Ablation (WMV+D-3):}
        Pruned 4-model subset achieves \textbf{88.4\%} (175/198) —
        the project-high accuracy.
  \item \textbf{Applied tasks (Phase~18):}
        5-model swarm matches Claude~Sonnet~4.5 on security audit
        (10/10 bugs) at comparable cost.
  \item \textbf{GovSim (Phase~19):}
        0\% $\to$ 100\% survival (Fisher $p=0.008$); CARE~$=0.97$,
        Gini~$=0.001$.
\end{itemize}

Total cost: under \$95 USD across $\sim$14,500 API calls.

% =====================================================================
\section{Introduction}\label{sec:intro}
% =====================================================================

\subsection{Problem Statement}

The proliferation of large language models from competing vendors has
created a fragmented ecosystem in which models cannot natively
coordinate.  Each architecture processes language through different
tokenization, attention mechanisms, training corpora, and alignment.
Yet many applications demand coordination: multi-agent reasoning,
ensemble decision-making, negotiation, and collaborative
problem-solving.

Recent work on multi-agent debate (MAD) has demonstrated that
structured interaction between LLMs can improve factuality and
reasoning~\cite{du2024debate}.  However, the MAD paradigm faces
critical challenges: unstructured debate can \emph{degrade} reasoning
quality~\cite{smit2024mad}; strong single-agent prompting can match
multi-agent systems on many benchmarks~\cite{wang2024rethinking}; and
existing MAS-LLM systems generally lack formal foundations, statistical
rigor, and honest engagement with failure modes~\cite{lamalfa2025limits}.

The \arrival Protocol addresses these challenges through three
complementary innovations:

\begin{enumerate}
  \item \textbf{Structured semantic atoms}: 66 compositional @-atoms
    provide a formal coordination vocabulary, preventing quality
    degradation.
  \item \textbf{CRDT-based conflict resolution}: MEANING-CRDT v1.1
    formalizes consensus as weighted vector averaging with proven
    convergence, bounded debt, and Bayesian equivalence.
  \item \textbf{Empirical rigor}: 2,200+ experiments across 17
    architectures with significance testing, adversarial testing,
    ablation, and engagement with critical work.
\end{enumerate}

\subsection{Contributions}

\begin{itemize}
  \item \textbf{\arrival Protocol}: 4-round structured coordination
    using 66 semantic atoms, validated across 17 LLM architectures.
  \item \textbf{MEANING-CRDT v1.1}: Mathematical framework with
    11 theorems for semantic conflict resolution.
  \item \textbf{Echo-chamber analysis}: 7 quantitative metrics for
    detecting social conformity in LLM ensembles.
  \item \textbf{Hypercorrection paradox}: Naive memory degrades
    multi-agent accuracy; gated CARE-ALERT restores it.
  \item \textbf{Applied validation}: 5-model swarm matches frontier
    solo at comparable cost.
  \item \textbf{CARE as task discriminator}: Quantitatively
    distinguishes analytical (CARE~$\approx0.50$) from constructive
    (CARE~$\approx0.97$) tasks.
  \item \textbf{Cooperation breakthrough}: First structured
    multi-agent protocol solving tragedy of the commons for LLMs
    (0\%$\to$100\% survival in GovSim).
  \item \textbf{CGD}: Agreement-gated debate achieving 86.4\% on GPQA
    Diamond with mid-tier models.
  \item \textbf{Non-monotonic scaling}: Quality~$>$~quantity —
    pruned 4-model subset outperforms 7-model ensemble.
  \item \textbf{Comprehensive evaluation}: 2,200+ experiments, under
    \$95, with 34 references to related work.
\end{itemize}

% =====================================================================
\section{Related Work}\label{sec:related}
% =====================================================================

\subsection{Multi-Agent Debate}

Du et al.~\cite{du2024debate} established that multi-agent debate
improves factuality, founding the MAD paradigm.  Liang et
al.~\cite{liang2024divergent} introduced structured critique rounds.
Li et al.~\cite{li2024sparse} showed sparse communication topologies
outperform fully connected ones.  RUMAD~\cite{wang2026rumad} unifies
debate with reinforcement learning.  A-HMAD~\cite{zhou2025ahmad}
proposes adaptive heterogeneous debate.  ReConcile~\cite{chen2024reconcile}
demonstrates round-table consensus among diverse LLMs.

MoA~\cite{togetherai2024moa} demonstrated that layered aggregation
outperforms individual models.  Li et al.~\cite{li2025selfmoa} showed
homogeneous ensembles can outperform heterogeneous ones — confirmed by
our Phase~16.  Hegazy~\cite{hegazy2024smaller} examines when smaller
models debate larger ones.  Ashery et al.~\cite{ashery2025social} study
social convention formation in LLM populations.

\subsection{Critical Perspectives}

Wang et al.~\cite{wang2024rethinking} showed single-agent CoT can
match multi-agent systems.  Smit et al.~\cite{smit2024mad} showed
debate can \emph{degrade} reasoning.  Wu et al.~\cite{wu2025debate}
question whether LLM agents can really debate in a controlled logical
reasoning study.  La Malfa et
al.~\cite{lamalfa2025limits} critiqued MAS-LLM research for lacking
formal foundations.  Choi et al.~\cite{choi2025debate} proved a
\textbf{martingale property}: debate does not change the expected
correctness of LLM responses — validating our Phase~21 finding that
debate is net zero on MCQ accuracy.

\subsection{Game-Theoretic Cooperation}

GovSim~\cite{piatti2024govsim}: 12/15 LLM models achieve 0\% survival
in commons games; reproduced by Curvo et al.~\cite{curvo2025govsim}.
FAIRGAME~\cite{buscemi2025fairgame}: end-game
defection.  Pich\'{e} et al.~\cite{piche2025robust}: RL shifts toward
defection.  Kaesberg et al.~\cite{kaesberg2025defection} investigate
LLMs' failure to avoid tragedy of commons.
Akata et al.~\cite{akata2023repeated} study repeated games with LLMs.
SanctSim~\cite{sanctsim2025} simulates effects of sanctions.
Hintze \& Adami~\cite{hintze2026cooperation}: cooperation emerges at a
tipping point.  Kostka \& Chudziak~\cite{kostka2026tom} evaluate
theory of mind in multi-agent systems.

\subsection{Selective Debate and Confidence Gating}

Fan et al.~\cite{fan2026imad} propose \textbf{iMAD}: ML-trained
confidence classifiers to gate debate.  Eo et al.~\cite{eo2025down}
examine whether natural language is optimal for debate
(\textbf{DOWN}).  Kaesberg et al.~\cite{kaesberg2025voting}
systematically compare voting versus consensus.  Ai et
al.~\cite{ai2025majority} go beyond majority voting with higher-order
aggregation.  Our CGD uses zero-shot inter-model agreement — requiring
no training data.

% =====================================================================
\section{The ARRIVAL Protocol}\label{sec:protocol}
% =====================================================================

\subsection{Atom Taxonomy}

The protocol employs 66 structured semantic @-atoms from
DEUS.PROTOCOL~v0.5:

\begin{definitionbox}[Core Atoms (46)]
Identity (\texttt{@SELF}, \texttt{@OTHER}), Goals (\texttt{@GOAL}),
Reasoning (\texttt{@REASON}, \texttt{@EVIDENCE},
\texttt{@HYPOTHESIS}), Conflict (\texttt{@CONFLICT},
\texttt{@DISAGREE}), Consensus (\texttt{@CONSENSUS},
\texttt{@RESOLUTION}), Meta (\texttt{@META}, \texttt{@REFLECT}).
\end{definitionbox}

\begin{definitionbox}[Promoted Emergent Atoms (20)]
Discovered in Phase~4 experiments (5/5 architecture adoption,
frequency $\geq 5$): Planning (\texttt{@ACTION\_PLAN},
\texttt{@PROPOSAL}), Evaluation (\texttt{@EVALUATION},
\texttt{@SYNTHESIS}), Risk/Knowledge (\texttt{@RISK\_TOLERANCE},
\texttt{@UNCERTAINTY}).
\end{definitionbox}

\subsection{Four-Round Protocol}

\begin{enumerate}
  \item \textbf{R1 --- Independent Analysis:} Each agent independently
    analyzes the question using \texttt{@HYPOTHESIS},
    \texttt{@EVIDENCE}, \texttt{@CONFIDENCE}.
  \item \textbf{R2 --- Adversarial Peer Review:} Each agent reviews
    others' R1 responses using \texttt{@CONFLICT},
    \texttt{@COUNTER}, \texttt{@DISAGREE}.
  \item \textbf{R3 --- CRDT Overlay:} Zero-cost mathematical
    reconciliation: CARE Resolve computes weighted consensus.
  \item \textbf{R4 --- Synthesis:} Designated synthesizer produces
    final consensus answer.
\end{enumerate}

\subsection{Confidence-Gated Debate (CGD)}

Phase~22 introduces an agreement-gated variant:

\begin{resultbox}[CGD Protocol]
\begin{enumerate}
  \item \textbf{Independent Solo} — All $K$ models answer
    independently (no inter-model exposure).
  \item \textbf{Agreement Check} — Classify: unanimous ($K/K$),
    split 2v1, split 3-way.
  \item \textbf{Targeted Debate} — Only when disagreement exists;
    final answer by majority vote (no R4).
\end{enumerate}
\end{resultbox}

This eliminates anchoring bias and R4 regression problems.

% =====================================================================
\section{MEANING-CRDT v1.1}\label{sec:crdt}
% =====================================================================

\noindent
We formalize conflict resolution through MEANING-CRDT
v1.1~\cite{kelevra2026crdt}, a mathematical framework based on
Conflict-free Replicated Data Types~\cite{shapiro2011crdt}.

\subsection{CARE Resolve}

\begin{theorembox}[Theorem 1 (Optimality)]
CARE Resolve minimizes total weighted semantic distance:
\[
  \hat{\mathbf{v}} = \arg\min_{\mathbf{v}}
    \sum_{i=1}^{N} w_i \, d(\mathbf{v}, \mathbf{v}_i)^2
    = \frac{\sum_{i} w_i \mathbf{v}_i}{\sum_{i} w_i}
\]
\end{theorembox}

\begin{theorembox}[Theorem 3 (Factor-4 Improvement)]
Under mild conditions, CARE provides up to $4\times$ reduction in
total semantic distance compared to unweighted averaging.
\end{theorembox}

\begin{theorembox}[Theorem 5 (Convergence)]
Iterated CARE converges to a fixed point in $O(N)$ rounds for $N$
agents with bounded weight ratios.
\end{theorembox}

\subsection{Meaning Debt}

\begin{theorembox}[Theorem 6 (Bounded Debt)]
Under the \arrival protocol, Meaning Debt is bounded:
$\mathrm{MD} \leq w_{\max} \cdot D_{\max} \cdot N$.
\end{theorembox}

\begin{theorembox}[Theorem 8 (Bayesian Equivalence)]
CARE Resolve is equivalent to Bayesian belief updating when confidence
weights correspond to prior precision parameters.
\end{theorembox}

\subsection{Vulnerability Analysis}

\begin{theorembox}[Theorem 9 (Weight Inflation)]
An adversarial agent can exploit CARE by inflating its confidence.
The damage is bounded by $w_{\mathrm{adv}} / \sum w_i$, making CARE
robust when honest agents' total weight dominates.
\end{theorembox}

% =====================================================================
\section{Experimental Setup}\label{sec:setup}
% =====================================================================

\subsection{Models and Infrastructure}

We evaluate across 17 LLM architectures (Table~\ref{tab:models}).

\begin{table}[H]
\centering\small
\caption{Models used across 23 experimental phases.}\label{tab:models}
\begin{tabular}{@{}llll@{}}
\toprule
\textbf{Model} & \textbf{Provider} & \textbf{Params} & \textbf{Phases} \\
\midrule
GPT-4o           & OpenAI      & Undisclosed  & 4--12 \\
Claude 3.5 Sonnet& Anthropic   & Undisclosed  & 4--12 \\
DeepSeek V3      & DeepSeek    & 671B MoE     & 4--13 \\
DeepSeek R1      & DeepSeek    & 671B MoE     & 4--13 \\
Llama 3.3 70B    & Meta        & 70B          & 4--12 \\
Qwen 2.5 72B     & Alibaba     & 72B          & 4--12 \\
Mistral Large    & Mistral     & Undisclosed  & 4--12 \\
Gemini 2.0 Flash & Google      & Undisclosed  & 4--12 \\
Qwen3-235B       & Alibaba     & 235B MoE     & 13, 16, 17 \\
Claude Sonnet 4.5& Anthropic   & Undisclosed  & 18 \\
GPT-4.1          & OpenAI      & Undisclosed  & 18, 20--22 \\
DeepSeek V3.2    & DeepSeek    & MoE          & 18, 23 \\
Mistral Large 3  & Mistral     & Undisclosed  & 18 \\
Gemini 3 Flash   & Google      & Undisclosed  & 18, 20--23 \\
Grok 4.1 Fast    & xAI         & Undisclosed  & 18, 20--23 \\
Claude Sonnet 4.6& Anthropic   & Undisclosed  & 23 \\
Kimi K2.5        & Moonshot AI & Undisclosed  & 23 \\
GLM-5            & Zhipu AI    & Undisclosed  & 23 \\
Qwen3.5-397B     & Alibaba     & 397B MoE     & 23 \\
\bottomrule
\end{tabular}
\end{table}

All models accessed via OpenRouter API.  No model fine-tuning; all
coordination via system prompt injection.

\subsection{Statistical Methods}\label{sec:stats}

All pairwise accuracy comparisons use \textbf{McNemar's test} (exact,
two-sided) for matched pairs on the same question set.
\textbf{Fisher's exact test} is used when comparing independent
conditions (e.g., GovSim survival rates).  Confidence intervals are
\textbf{Wilson score} intervals at 95\%.  Effect sizes are reported as
accuracy differences in percentage points (\pp).  Where applicable,
\textbf{Mann--Whitney~U} is used for non-parametric comparisons of
continuous metrics (CARE scores, Meaning Debt).  No correction for
multiple comparisons is applied; all $p$-values are reported for
transparency.

\subsection{Benchmark}

GPQA Diamond~\cite{rein2024gpqa}: graduate-level science, resistant
to web search.  Phases~13/16/17 use 40 questions; Phases~20--23 use
the full 198-question set (physics~86, chemistry~93, biology~19).

Current GPQA Diamond SOTA: 94.1\% (Gemini 3.1 Pro~\cite{gemini2025}).
Human expert accuracy: 69.7\%.
As of March 2026, multiple frontier models exceed 90\%: Gemini 3.1 Pro
94.1\%, GPT-5.3 Codex $\sim$91.5\%.
Kim et al.~\cite{kim2025scaling} investigate scaling laws in
multi-agent systems.

\subsection{Budget}

\begin{table}[H]
\centering\small
\caption{Cost breakdown across all 23 experimental phases.}\label{tab:budget}
\begin{tabular}{@{}lrrr@{}}
\toprule
\textbf{Phase} & \textbf{Experiments} & \textbf{API Calls} & \textbf{Cost (USD)} \\
\midrule
4--5   & 485     & $\sim$2,000 & $\sim$\$2.00 \\
6--12  & $\sim$80& $\sim$500   & $\sim$\$1.50 \\
13     & 80      & $\sim$320   & $\sim$\$1.00 \\
14--15 & $\sim$40& $\sim$200   & $\sim$\$0.50 \\
16     & 200+    & 640         & \$1.92 \\
17     & 200     & 200         & \$0.50 \\
18     & 6       & $\sim$66    & \$0.77 \\
19     & 25      & 2,165       & \$5.76 \\
20     & 198     & $\sim$1,782 & \$7.43 \\
21     & 792     & $\sim$1,386 & \$9.63 \\
22     & 198     & $\sim$692   & \$5.72 \\
23     & 198     & $\sim$1,822 & \$11.35 \\
AutoGen& 16      & $\sim$100   & \$0.02 \\
\midrule
\textbf{Total} & \textbf{2,200+} & \textbf{$\sim$14,500} & \textbf{$<$\$95} \\
\bottomrule
\end{tabular}
\end{table}

% =====================================================================
\section{Results: Core Protocol (Phases 4--12)}\label{sec:core}
% =====================================================================

In 385 dyadic coordination experiments across 8 LLM architectures,
cross-architecture consensus reached \textbf{98.6\%} (142/144 unique
pairings).  Models spontaneously generated 506 emergent atoms beyond
the standard 46, with 1,173 cross-architecture adoptions.

Adversarial testing (Phases~6--7): Byzantine saboteurs with
\texttt{@SABOTAGE} atoms were detected and neutralized by honest
agents in $>$90\% of cases.  Trojan atoms (designed to pass
validation) were identified and rejected within 2 rounds.

% =====================================================================
\section{Results: GPQA Diamond (Phase 13, N=40)}\label{sec:gpqa40}
% =====================================================================

\begin{resultbox}[Phase 13 Key Result]
\arrival 63.8\% vs.\ Majority Voting 42.5\%
(McNemar $p = 0.006$, significant on 40 questions).
\end{resultbox}

Per-domain: Biology +40\pp (80\% vs.\ 40\%), Chemistry +24\pp (57\%
vs.\ 33\%), Physics +5\pp (57\% vs.\ 52\%).

% =====================================================================
\section{Results: Homogeneous Ensemble (Phase 16)}
% =====================================================================

Phase~16 tests a 5-copy Qwen3-235B ensemble with different personas.
Result: \textbf{65.0\%} (26/40) vs.\ baseline MV 52.5\% (+12.5\pp,
McNemar $p=0.006$).  Echo-chamber analysis reveals mean uniqueness
ratio of 0.73 (sufficient diversity despite identical architecture).

% =====================================================================
\section{Results: Solo CoT Baseline (Phase 17)}
% =====================================================================

Solo Qwen3-235B with enhanced CoT: \textbf{70.0\%} MV accuracy,
5\pp above \arrival's 65.0\% (not significant, $p=0.812$).  Validates
Wang et al.'s finding that strong single-agent prompting can match
multi-agent on accuracy.

% =====================================================================
\section{Results: Full GPQA Diamond (Phases 20--23)}\label{sec:gpqa198}
% =====================================================================

\subsection{Phase 20: Mid-Tier Models (N=198)}

\begin{resultbox}[Phase 20 Key Result]
\arrival MV \textbf{67.7\%} vs.\ Solo CoT MV 62.1\%
(+5.6\pp, McNemar $p=0.233$, ns).
\end{resultbox}

Per-domain breakdown:

\begin{table}[H]
\centering\small
\caption{Phase 20 per-domain accuracy (N=198).}\label{tab:phase20domain}
\begin{tabular}{@{}lrrr@{}}
\toprule
\textbf{Domain} & \textbf{\arrival MV} & \textbf{Solo CoT MV} & \textbf{$\Delta$} \\
\midrule
Biology (19)   & 84.2\% & 63.2\% & +21.1\pp \\
Chemistry (93) & 64.5\% & 58.1\% & +6.5\pp \\
Physics (86)   & 67.4\% & 66.3\% & +1.2\pp \\
\midrule
\textbf{Overall (198)} & \textbf{67.7\%} & \textbf{62.1\%} & \textbf{+5.6\pp} \\
\bottomrule
\end{tabular}
\end{table}

\subsection{Phase 21: Frontier Models (N=198)}

\begin{resultbox}[Phase 21 Key Result --- Debate Effect Inverts]
Non-Debate MV \textbf{82.3\%} vs.\ \arrival MV 78.8\%
($-3.5$\pp, $p=0.265$, ns).\\
\textbf{Anchoring}: Gemini drops from 82.3\% (solo) to 78.3\% (debate)
--- a 4.0\pp degradation.
\end{resultbox}

R4 finalizer is consistently net negative: Phase~20 (9 regressions vs.\
7 rescues), Phase~21 (13 vs.\ 8).

\begin{table}[H]
\centering\small
\caption{Phase 21 voting strategy comparison (N=198).}\label{tab:phase21voting}
\begin{tabular}{@{}lr@{}}
\toprule
\textbf{Strategy} & \textbf{Accuracy} \\
\midrule
Oracle (any of 3 correct) & 93.9\% \\
Non-Debate MV             & 82.3\% \\
Grok-weighted MV          & 86.4\% \\
Solo Grok 4.1 Fast        & 85.4\% \\
Solo GPT-4.1              & 78.3\% \\
Solo Gemini 3 Flash       & 76.3\% \\
\arrival R4               & 76.3\% \\
\arrival MV               & 78.8\% \\
\bottomrule
\end{tabular}
\end{table}

\subsection{Phase 22: Confidence-Gated Debate (N=198)}

\begin{resultbox}[Phase 22 Key Result --- CGD]
\cgd achieves \textbf{86.4\%} accuracy (171/198).\\
95\% CI: [80.9\%, 90.5\%].\\
vs.\ \arrival MV: +7.6\pp (McNemar $p=0.009$, significant).\\
61.6\% questions resolved unanimously (95.9\% accuracy).\\
Cost: \$5.72 for 198 questions (\$0.029/q).
\end{resultbox}

\begin{table}[H]
\centering\small
\caption{Phase 22 CGD results on GPQA Diamond (198 questions).}\label{tab:cgd3}
\begin{tabular}{@{}lrrl@{}}
\toprule
\textbf{Method} & \textbf{Correct} & \textbf{Accuracy} & \textbf{Notes} \\
\midrule
\textbf{CGD (Phase 22)}   & 171/198 & \textbf{86.4\%} & Primary result \\
Solo Grok 4.1 Fast        & 169/198 & 85.4\%           & Best solo \\
Non-Debate MV             & 163/198 & 82.3\%           & Phase 21 \\
\arrival MV               & 156/198 & 78.8\%           & Phase 21 \\
Oracle (any of 3)         & 189/198 & 95.5\%           & Diversity ceiling \\
\bottomrule
\end{tabular}
\end{table}

\begin{table}[H]
\centering\small
\caption{Phase 22 CGD per-domain accuracy.}\label{tab:cgddomain}
\begin{tabular}{@{}lrr@{}}
\toprule
\textbf{Domain} & \textbf{Correct/Total} & \textbf{Accuracy} \\
\midrule
Physics (86)   & 82/86  & \textbf{95.3\%} \\
Chemistry (93) & 74/93  & 79.6\% \\
Biology (19)   & 15/19  & 78.9\% \\
\midrule
\textbf{Overall (198)} & \textbf{171/198} & \textbf{86.4\%} \\
\bottomrule
\end{tabular}
\end{table}

\textbf{Minority-Was-Right}: In 65 split 2v1 questions, minority
correct 15 times (23.1\%).  GPT-4.1 as minority: 1/32 = 3\%.
Grok: 8/18 = 44\%.

\subsection{Phase 23: 7-Model Cross-Vendor CGD}

\begin{table}[H]
\centering\small
\caption{Phase 23 model lineup: 3 US + 4 Chinese models from 5 vendors.}\label{tab:phase23models}
\begin{tabular}{@{}clllr@{}}
\toprule
\# & \textbf{Model} & \textbf{Vendor} & \textbf{Country} & \textbf{Solo} \\
\midrule
1 & Grok 4.1 Fast   & xAI         & US & \textbf{87.4\%} \\
2 & Gemini 3 Flash   & Google      & US & 82.8\% \\
3 & Claude Sonnet 4.6& Anthropic   & US & 80.8\% \\
4 & DeepSeek V3.2    & DeepSeek    & CN & 74.7\% \\
5 & Kimi K2.5        & Moonshot AI & CN & 72.2\% \\
6 & GLM-5            & Zhipu AI    & CN & 70.7\% \\
7 & Qwen3.5-397B     & Alibaba     & CN & 58.6\% \\
\bottomrule
\end{tabular}
\end{table}

\begin{resultbox}[Phase 23 Key Result --- CGD-7]
CGD-7: \textbf{86.9\%} (172/198) vs.\ Phase 22 CGD: 86.4\%
(+0.5\pp, $p=1.0$, ns).\\
Best solo (Grok): 87.4\% --- solo $>$ ensemble.\\
Oracle ceiling: 96.5\% (191/198).\\
Cost: \$11.35 total.
\end{resultbox}

\subsubsection{Debate Type Analysis}

\begin{table}[H]
\centering\small
\caption{CGD-7 debate type breakdown (198 questions).}\label{tab:debate}
\begin{tabular}{@{}lrrl@{}}
\toprule
\textbf{Type} & \textbf{N (\%)} & \textbf{Accuracy} & \textbf{Action} \\
\midrule
Unanimous (7/7)    & 105 (53\%) & 99.0\% & Locked \\
Supermajority (6/7)& 23 (12\%)  & 82.6\% & Locked \\
Strong maj.\ (5/7) & 22 (11\%)  & 81.8\% & Locked \\
Simple maj.\ (4v3) & 21 (11\%)  & 66.7\% & Minority debate \\
No majority        & 27 (14\%)  & 63.0\% & Full debate \\
\bottomrule
\end{tabular}
\end{table}

75.8\% of questions locked without debate ($\geq$5/7 agreement),
achieving 93.3\% accuracy.

\subsubsection{Ablation: Weighting, Pruning, Poisoned Voters}

\textbf{Weighted vs.\ Simple}: Weighting \emph{hurts} — Weighted MV
(85.4\%) $<$ Simple MV (85.9\%) $<$ CGD-7 (86.9\%).

\textbf{WMV+D-3 (pruned 4-model subset)}: Dropping the 3 weakest
Chinese models yields \textbf{88.4\%} (175/198) — the project-high
accuracy.

\textbf{Poisoned voter}: Qwen3.5-397B exhibits 5\% minority-correct
rate with 54 extraction failures (27\%).  Its dissent reliably
indicates the \emph{wrong} answer.

\textbf{Thinking mode confound}: US models averaged 83.7\% vs.\ CN
69.1\% — but reasoning was disabled for all Chinese models, making
this comparison confounded.

% =====================================================================
\section{Results: Applied Tasks (Phase 18)}\label{sec:applied}
% =====================================================================

\begin{resultbox}[Phase 18: Five Models Match One Frontier Model]
5-model \arrival swarm matches Claude Sonnet 4.5:\\
Security audit: 10/10 bugs found (both).\\
Code generation: 53\% vs.\ 60\% test pass rate.\\
Cost: \$0.14 vs.\ \$0.12 per task.
\end{resultbox}

\care as task discriminator: analytical tasks (security audit)
produce CARE~$\approx0.50$; constructive tasks (code generation)
produce CARE~$\approx0.97$.

% =====================================================================
\section{Results: Memory and CARE-ALERT (Phases 14--15)}
% =====================================================================

\textbf{Hypercorrection paradox (Phase~14)}: Naive memory injection
degrades accuracy by 5.7\pp.

\textbf{Gated CARE-ALERT (Phase~15)}: Monitors Meaning Debt in
real time; restores baseline performance (CARE improvement $p=0.042$).
Validated in Phase~18: correctly activates when CARE~$=0.50$ and
suppresses when CARE~$=0.97$.

% =====================================================================
\section{Results: AutoGen Validation}\label{sec:autogen}
% =====================================================================

Framework-agnostic validation on Microsoft AutoGen
(AG2)~\cite{wu2023autogen} achieves
100\% behavioral match with the reference implementation across
16 test cases.

% =====================================================================
\section{Results: GovSim (Phase 19)}\label{sec:govsim}
% =====================================================================

\begin{resultbox}[Phase 19: Cooperation Breakthrough]
\arrival converts \textbf{0\% baseline survival to 100\% survival}
across all 25 runs (5 conditions $\times$ 5 runs).\\
Fisher exact $p = 0.008$.\\
CARE $= 0.97$, Gini $= 0.001$ (near-perfect equality).\\
$3.2\times$ resource extraction vs.\ baseline collapse.
\end{resultbox}

\textbf{Key findings}:
\begin{enumerate}
  \item \textbf{Self-healing}: Greedy proposals neutralized within
    1 round via CARE convergence.
  \item \textbf{Condition D (binding-only)}: 100\% survival at
    $32\times$ lower cost.
  \item \textbf{Condition E (adversarial R4)}: Survival preserved
    (safety clamp) but fairness destroyed (Gini $= 0.800$).
  \item \textbf{Empirical support for Theorems 1, 5, 6, 8}.
\end{enumerate}

% =====================================================================
\clearpage
\section{Discussion}\label{sec:discussion}
% =====================================================================

\subsection{What ARRIVAL Adds Beyond Strong Prompting}

\arrival does not claim to always improve raw accuracy.  Its value
lies in: (1)~structured reasoning traces, (2)~error detection via
peer review, (3)~formal guarantees via MEANING-CRDT, (4)~scalability,
and (5)~protocol evolution — CGD emerged from Phase~20--21 failure
analysis, achieving 86.4\%.

\subsection{Cost-Efficiency}

At under \$95 for $\sim$14,500 API calls across 23 phases and 17
architectures, \arrival demonstrates exceptional cost-efficiency.
Phase~22 CGD: \$5.72 for 198 questions (\$0.029/q).
Phase~23 CGD-7: \$11.35 for 198 questions (\$0.057/q).
Multi-agent overhead: $\sim$\$0.012--\$0.057 per question.

\subsection{CGD: From Failure Analysis to Protocol Innovation}

CGD addresses three Phase~20--21 problems: (1)~anchoring from
sequential exposure, (2)~destructive R4 overrides, (3)~false
consensus.  Grounded in the Condorcet Jury Theorem: independent
voters with accuracy $>50\%$ produce better collective decisions.
Choi et al.~\cite{choi2025debate} prove that debate has a martingale
property — CGD sidesteps this by reserving debate for genuine
disagreement.

\subsection{Quality Threshold and Non-Monotonic Scaling}

Phase~23 reveals non-monotonic scaling: 3$\to$7 models yields only
+0.5\pp ($p=1.0$), while pruned 4-model WMV+D-3 achieves 88.4\%.
Qwen3.5-397B: 5\% minority-correct, 27\% extraction failures —
a \textbf{poisoned voter}.  Implication: model selection $>$ model
count.

\subsection{Thinking Mode as Hidden Variable}

US--CN accuracy gap (83.7\% vs.\ 69.1\%) is confounded by systematic
disabling of reasoning modes for Chinese models.  The 14.6\pp gap
should \textbf{not} be interpreted as inherent architectural
differences.

% =====================================================================
\section{Limitations}\label{sec:limitations}
% =====================================================================

\begin{enumerate}
  \item \textbf{Benchmark specificity}: GPQA Diamond results may not
    generalize to all reasoning tasks.
  \item \textbf{API dependence}: All experiments use API-accessible
    models; local/fine-tuned models untested.
  \item \textbf{Sequential bias}: The 4-round structure can introduce
    anchoring (addressed by CGD).
  \item \textbf{Cost of debate}: Multi-agent overhead may not be
    justified when a single frontier model suffices.
  \item \textbf{CARE sensitivity}: Heuristic weight extraction from
    confidence declarations.
  \item \textbf{No social agency}: LLM agents lack genuine social
    agency~\cite{lamalfa2025limits}.
  \item \textbf{Homogeneous caveat}: Phase~16 uses 5 copies of the
    same model.
  \item \textbf{Phase 18 N=1}: Applied experiments lack repeated
    trials.
  \item \textbf{Code synthesis}: 4-round protocol incurs coherence
    loss.
  \item \textbf{Single game type}: GovSim fish pond only.
  \item \textbf{No adversarial agents}: Cooperative personas only.
  \item \textbf{CGD limited scope}: Tested on MCQ with $K=3$ and
    $K=7$; non-MCQ untested.
\end{enumerate}

% =====================================================================
\section{Conclusion and Future Work}\label{sec:conclusion}
% =====================================================================

We presented the \arrival Protocol, demonstrating that structured
semantic communication enables multi-agent LLM coordination across
benchmarks, practical tasks, and cooperation scenarios.  Key results:

\begin{itemize}
  \item \textbf{Phase 22 CGD}: 86.4\% on GPQA Diamond (171/198),
    +7.6\pp vs.\ \arrival MV ($p=0.009$).
  \item \textbf{Phase 23 CGD-7}: 86.9\% with 7 models; non-monotonic
    scaling confirmed.
  \item \textbf{WMV+D-3}: 88.4\% (175/198) — project-high accuracy.
  \item \textbf{Phase 19}: 0\%$\to$100\% survival in GovSim
    ($p=0.008$).
  \item \textbf{Phase 18}: 5 cheaper models match frontier solo.
\end{itemize}

\textbf{Honest picture}: standard \arrival does not significantly
outperform solo models on MCQ at scale; with frontier models, debate
hurts via anchoring.  CGD demonstrates that failure analysis was
productive: 86.4\% is the highest in the project.  \arrival's primary
value: cooperation (GovSim), transparency, vendor independence, and
principled ensemble aggregation.

\subsection*{Future Work}

\begin{enumerate}
  \item Extend CGD to additional benchmarks (MMLU-Pro, ARC-Challenge,
    MedQA).
  \item CGD extensions: $K > 3$ models, weighted confidence gating.
  \item Broader cooperation games: public goods, prisoner's dilemma.
  \item Adversarial agents: RL-trained defectors.
  \item Adaptive orchestration via real-time CARE measurements.
  \item Scaling laws: accuracy vs.\ number of agents.
  \item Online memory: cross-session coordination patterns.
  \item Larger N for MCQ ($N \approx 600$) for statistical power.
  \item CGD vs.\ frontier SOTA: cost-accuracy tradeoffs.
\end{enumerate}

% =====================================================================
% Bibliography
% =====================================================================
\bibliographystyle{plainnat}
\bibliography{references}

% =====================================================================
\appendix
\section{DEUS.PROTOCOL v0.5 Atom Dictionary}

See \texttt{docs/ATOM\_DICTIONARY.md} in the repository for the
complete specification of all 66 atoms with usage examples and
adoption statistics.

\section{MEANING-CRDT v1.1 Theorem Proofs}

Full proofs for all theorems are available at DOI:
\href{https://doi.org/10.5281/zenodo.18702383}{10.5281/zenodo.18702383}

\section{Experiment Logs and Reproducibility}

All experiment configurations, API logs, and result files are available
in the \texttt{experiments/} and \texttt{results/} directories.
Total cost: under \$95 USD across 2,200+ experiments and
$\sim$14,500 API calls.
See \texttt{docs/setup/SETUP\_EN.md} for reproduction instructions.

\end{document}
